\section{Multimodal Measurements}
\label{sec:methods}
Each task applies acquisition-quality and completion gates before interpreting measurements. For sampled tasks, $q_{\mathrm{valid}}=N_{\mathrm{valid}}/N_{\mathrm{captured}}$ uses counts within the component-specific analysis window. Missing landmarks and incomplete actions produce unavailable measurements, with explicit masks in the research representation. The operating points below define the engineering measurement protocol.

\subsection{Balance: Postural Stability}
\label{subsec:balance}
Passive B distinguishes stable standing from ordinary transitions such as
sitting, walking, or bending before evaluating change. A separate descent and
posture-change path can retain a fall-like trigger. Active B asks the user to
stand naturally for 30~s after a 1.5-s warmup. A qualified window requires
at least 30 valid samples and valid fraction at least 0.75. Shoulder, hip, and
ankle geometry supplies trunk
orientation, mediolateral motion, and central sway range. An active balance symptom exempts an unstarted standing attempt. A retry starts a new attempt; a failed earlier attempt remains incomplete in its historical report. Starting a standing task is blocked while a balance symptom remains active, including after an ordinary negative questionnaire answer. Only an explicit correction retracts that event.

Personal enrollment requires a symptom-free attestation and binds the baseline
to subject, device, and camera identifiers. The first five qualified windows
freeze the median $\mu_j$ and scale
\begin{equation}
\sigma_j=\max\{1.4826\operatorname{MAD}(X_j),
\operatorname{IQR}(X_j)/1.349\}.
\end{equation}
Changing the enrollment context clears its samples. For trunk orientation $\theta$, lateral-motion
velocity $v$, and central 90\% range $R$, define
\begin{align}
r_\theta&=\frac{|\theta-\mu_\theta|}{\sigma_\theta},\nonumber\\
r_v&=\frac{\max(0,v-\mu_v)}{\sigma_v},\quad
r_R=\frac{\max(0,R-\mu_R)}{\sigma_R}.
\end{align}
The rule is $r_\theta\geq3.5$ or both $r_v\geq3.5$ and $r_R\geq3.5$,
using a modified-$Z$ reference \cite{iglewicz1993outliers}. Requiring both sway
terms is intended to reduce landmark-jitter triggers. Zero change receives score zero, including at zero scale. Positive change at zero scale is unscorable. A scored positive rule takes precedence; otherwise, an unscorable domain yields insufficient evidence. Figure~\ref{fig:balance} shows acquisition before the
personal baseline has been completed.



\subsection{Eyes: Gaze Assessment}
The E task repeats paired center--left and center--right target trials three
times. Each target lasts 2~s, with the first 0.5~s reserved for settling. Normalized iris
positions produce paired endpoint responses, range, directional asymmetry, and
inter-eye differences. Each trial requires five valid samples and valid fraction at least 0.60; each eye must span at least 24 pixels when frame dimensions are available. Repeated responses, head pose, and endpoint dispersion further gate availability. The accompanying measurement protocol fixes these parameters and their definitions. These RGB outputs are
research measurements: automatic E disease thresholds are disabled. A separate
questionnaire records visual difficulty and its onset, even when camera
acquisition fails.

\subsection{Face: Facial Asymmetry}
F compares a 2-s neutral phase with a 3-s smile phase, requiring at least five
valid samples in each. Mouth-corner displacement is normalized
by interocular width; its median change is interpreted alongside left--right
smile activation. The trial requires valid fraction at least 0.50 and smile
strength at least 0.22. Qualified asymmetry is positive when the absolute
neutral-referenced geometric change reaches 0.075 or smile-activation difference
reaches 0.24. Failure to establish the action produces insufficient evidence;
users or caregivers can separately report inability to smile.

\subsection{Arms: Arm Drift}
\label{subsec:arms}
The A task follows a lower--raise--hold--return sequence, with a 5-s hold near
shoulder height (Fig.~\ref{fig:arms}). Bilateral visibility and phase completion
are required, with at least 20 visible hold samples and hold-phase valid fraction at least
0.55. The raised-pose gate applies when entering the hold phase. During the
hold, all frames with measurable shoulder, elbow, and wrist landmarks are
retained, including lowered-arm poses. Early and late segments occupy the first
and last thirds of the 5-s interval, with at least six visible samples in each.
An unavailable endpoint produces insufficient evidence. With image $y$ increasing downward, shoulder-width-normalized
wrist height is $h_{e,t}=(y_{wrist,e,t}-y_{shoulder,e,t})/w_{shoulder}$.
We separate persistent asymmetry from progressive lowering:
\begin{align}
D&=\operatorname{median}_t(h_{L,t}-h_{R,t}),\\
\delta_e&=\operatorname{median}(h_e^{late})-
\operatorname{median}(h_e^{early}),\quad
\Delta=\delta_L-\delta_R.
\end{align}
A qualified trial is positive when $|D|\geq0.30$ or $|\Delta|\geq0.22$.
Anatomical labels assume an unmirrored inference frame; positive differences
indicate a lower left arm. Technical inability to obtain this trial remains
separate from a reported inability to raise an arm.

\begin{figure}[t]
\centering
\includegraphics[width=0.80\columnwidth]{arms-demonstration.jpg}
\caption{Bilateral arm-raise demonstration with live landmark feedback on the local computer.}
\label{fig:arms}
\end{figure}

\begin{figure}[t]
\centering
\includegraphics[width=0.80\columnwidth]{balance-demonstration.jpg}
\caption{Balance acquisition with landmark feedback; baseline-relative scores remain unavailable during enrollment (0/5 windows).}
\label{fig:balance}
\end{figure}

\subsection{Speech: Audio Assessment}
Passive S compares qualified audio windows with an enrolled personal baseline.
The default loop acquires an 8-s window, analyzes it, then waits 1~s; this is
sequential window processing. Three abnormal votes are required among the last
five qualified windows. From an empty vote history at a window boundary, three
8-s captures and two 1-s waits require at least 26~s plus processing.
A target-speaker gate uses three enrollment windows and a cosine-similarity
reference of 0.70. The classifier probability contributes only one anomaly domain; a
trigger additionally requires another changed acoustic domain. The diagnostic
output links votes to window identifiers, capture boundaries, and analysis and
decision times for subsequent timing audits. The passive gate and the active check share classifier evidence; the active
check adds a new prompted recording and onset information.

Active S acquires a new 7-s, 16-kHz mono recording of a fixed Mandarin sentence.
After duration, level, clipping, and voice-content gates, a frozen logistic
classifier trained on MDSC supplies dysarthria probability. A value at or above 0.735282
is the direct S positive criterion. Local Whisper transcription supplies
character error rate (CER), speaking rate, and pause descriptors for audit; these
features do not override the classifier. Dysarthria is not stroke-specific,
so measurement onset and independent symptom reports remain necessary inputs
to the urgency rule.
